\section{怎样掌握文言的特殊句法}
句法是指组词成句的规则。由于句法结构的稳固性强，所以古今汉语的句法基本相同。但是语言是发展的，古今汉语的句法仍有不同的地方。这些不同的地方，就是我们学习文言文的困难之处。我们只要弄清这些不同的地方，就可以掌握文言的特殊句法。这些不同的地方简述如下：

1. 判断句的特殊表示法

在古代汉语中，判断句一般不是由现代汉语判断词“是”字来表示的，而是用“······者，······也”的格式来表示；有的文言判断句则仅在主语后面用“者”字提顿，或用“也”字煞句，甚至“者”，“也”都不用，构成“······，······”的格式。在译成现代汉语时，一律用“······是······”的译法来表示，“是”之前翻译主语部分，“是”之后翻译谓语部分。例如：

①诸葛孔明者，卧龙也。（《隆中对》）

②廉颇者，赵之良将也。（《廉颇蔺相如列传》）

③南阳刘子骥，高尚士也。（《桃花源记》）

④其巫，老女子也。（《西门豹治邺》）

⑤此劲敌也。（《冯婉贞》）

⑥此天子之气也。（《鸿门宴》）

⑦裴子者，裴秀。（《马钧传》）

⑧粟者，民之所种。（《论贵粟疏》）

⑨此用武之国。（《隆中对》）

⑩刘备天下枭雄。（《赤壁之战》）

除此之外，有时用副词“乃、则、即、皆、素”，动词“为”等表示判断。“为”可以译为“是”；“乃”可以译为“是，原来是”；“则、即”可以译为“就是”；“皆”可以译为“都是”；“素”可以译为“一向是”。还可以用副词“非”字表示否定性判断，相当于现代的“不是”。例如：

①项燕为楚将。（《陈涉世家》）

②臣乃市井鼓刀屠者。（《信陵君窃符救赵》）

③此则岳阳楼之大观也。（《岳阳楼记》）

④梁父即楚将项燕。（《史记·项羽本纪》）

⑤夫六国与秦皆诸侯。（《六国论》）

⑥且相如素贱人······（《廉颇蔺相如列传》）

⑦人非生而知之者，孰能无惑？（《师说》）

⑧予本非文人画士。（《病梅馆记》）

古汉语中也偶而有用“是”字表示判断的，如《桃花源记》中“问今是何世”。《木兰诗》中“不知木兰是女郎”，《卖炭翁》中“翩翩两骑来是谁？”等。但是，我们必须注意：千万不要把作主语用的指示代词“是”误认为判断词。例如：

①不知军之不可以退而谓之退，是谓縻军。（《谋攻》）

（不了解军队不能够后退，却命令军队后退，这叫做束缚军队。）

②是叶公非好龙也。（《叶公好龙》）

（这说明叶公并不是喜爱真龙啊。）

被动句的特殊表示法

在古代汉语中，除了偶而用“被”字表示被动之外，一般是用“为······所······”、“见······于······”、“受······于······”等固定格式来表示，或者是用“见”、“为”、“于”等词来表示。例如：

①赢兵为人马所蹈藉，陷泥中。（《赤壁之战》）

（老弱残兵被人马践踏，陷到泥里边。）

②臣诚恐见欺于王而负赵。（《廉颇蔺相如列传》）（我实在耽心被大王欺骗而辜负了赵国。）

③吾不能举全吴之地，十万之众，受制于人。（《赤壁之战》）

（我不能总合整个吴国的土地，十万的兵士，被别人挟制。）

④欲予秦，秦城恐不可得，徒见欺。（《廉蔺列传》） （想要把和氏璧给秦国，只怕秦国的城池不能得到，白白地被欺骗。）

⑤吾属今为之虏矣。（《鸿门宴》）

（我们如今将被沛公俘虏啦。）

⑥暴见于王，王语暴以好乐，暴未有以对也。（《孟子二章·庄暴见孟子》）

（我庄暴被大王召见，大王把爱好音乐的事告诉我，我没有什么话拿来回答。）

省略句的一般规律

在一定的语言环境里，一些词语或者承前省略了，或者蒙后省略了，这种词语省略的现象是古今汉语所共有的。但是，文言里这种现象更为突出，甚至在现代汉语中一般不可省略的句子成分，文言里也经常被省略。对于这后一种情况，我们在学习文言文时，必须注意辨析。例如：

①永州之野产异蛇，（ ）黑质而白章，（ ）触草木，（ ）尽死。（《捕蛇者说》）

（永州的郊外出产一种奇异的蛇，这种蛇黑底子、白花纹，它碰上草木，草木就得死。）

②一鼓作气，再（ ）而（ ）衰，三（ ）而（ ）竭。（《曹刿论战》）（第一次击鼓能激发士气，第二次击鼓士气就开始衰落，第三次击鼓士气则已消失。）

③（ ）曰：“公之视廉将军孰与秦王（ ）？”（ ）曰：“（ ）不若（ ）（ ）也。（《廉颇蔺相如列传》）(蔺相如问道：“诸位看廉将军跟秦王相比哪个厉害？”大家回答说：“廉将军不如秦王厉害”。)（按：“孰与秦王”之后，省略了“威”字，因为下句就说“夫以秦王之威”，这是蒙后省略。）

④衣食所安，（ ）弗敢专也，必以（ ）分人。（《曹刿论战》）

（衣食这些使人安全的东西，我不敢独自享受，一定把它们分给别人。）

⑤扶苏以数谏故，上使（ ）（ ）外将兵。（《陈涉世家》）（扶苏由于屡次进谏的缘故，皇上派他到外边带兵。）

⑥将军战（ ）河（ ）北，臣战（ ）河（ ）南。（《鸿门宴》）

（将军您在黄河的北面作战，我在黄河的南面作战。）

⑦陈胜、吴广皆（ ）次当行。（《陈涉世家》）（陈胜、吴广按编次都应当应征出发。）

一般说来，古今汉语都可以承前或蒙后省略主语，或在对话的语言环境下省略主语（例如例句③），而文言文还可以在前后句主语不同的情况下省略主语，这需要联系上下文意，细心辨析，才能在译句时正确地补上所省略的主语，例如例句①。

省略谓语的情况，古今汉语都较少，但在文言文的并列句中，如果其中一句用了动词，另一句中同样的动词就经常省略，例如例句②。

在现代汉语里，宾语（尤其是介词的宾语）一般不能省略；在古代汉语里却都经常省略，例如例句③的动词“若”后面省略了宾语“之”（代“秦王”），例句④的介词“以”后面省略了宾语“之”（代“衣食”）。

在现代汉语里，兼语一般不能省略，古代汉语里的兼语却往往省略，例如例句⑤的使动词“使”之后省略了兼语“之”（代“扶苏”），因为这个“之”既是动词“使”的宾语，又是动词“将”的主语。

至于古代汉语中介词“于”的省略更为经常，例如例句⑥，有时介词“以”也可以省略，例如例句⑦。

文言的几种固定格式

（1）何如、如何、若何、奈何：相当于现代的“怎么”、“怎么样”或“怎么办”，不能译为“象什么”。例如：

①今日之事何如？（怎么样）（《鸿门宴》）

②取吾璧，不予我城，奈何？（怎么办）（《廉颇蔺相如列传》）

③民不畏死，奈何以死惧之！（怎么）（《老子》）

④使归就戮于秦，以逞寡君之志，若何？（怎么样）（《殽之战》）

（2）如······何、奈······何、若······何：相当于现代的“对（把）······怎么样（怎么办）”。例如：

①如太行、王屋何？（能把太行、王屋两座大山怎么样呢？）（《愚公移山》）

②其如土石何？（对这些土石又能怎么样呢？）（《愚公移山》）

③若楚惠何？（对楚国的恩惠怎么办？）（《左传·晋楚城濮之战》）

（3）如之何、奈之何、若之何。三个字作为固定结构，用作句子的谓语或状语，讲法就等于“如何、奈何、若何”，如果“之”是作为代词临时插进去的，讲法就等于“如······何”、“奈······何”、“若······何”。例如：

①巫妪、三老不来还，奈之何？（“对这事怎么办？” 这儿的“之”是作为代词临时插进去的。） （《西门豹治邺》）

②无鸡者，弗食鸡则已耳，去饥寒犹远，若之何去夫猫也？（“怎么赶走那猫呢？”这儿的“若之何”作为固定结构，用作句子的状语。） （刘基《郁离子》）

（4）孰与、孰若：相当于“同······比谁······”、“哪如······”。例如：

①我孰与城北徐公美？（我同城北徐公比谁美呢？）（《邹忌讽齐王纳谏》）

②为两郎僮，孰若为一郎僮耶？（做两个人的僮仆，哪如做一个人的僮仆呢？）（《童区寄传》）

（5）与其······孰若······：相当于“与其······，还不如（哪如）······”例如：

①与其坐且待亡，孰若起而拯之！（《冯婉贞》）

②与其杀是僮，孰若卖之？（《童区寄传》）

（6）无乃······乎、其······乎、得无······乎：相当于“恐怕（大概）······吧”，“能不（莫不是）······吗（吧）”。

例如：

① 师劳力竭，远主备之，无乃不可乎（恐怕不行吧）？（《殽之战》）

②王之好乐甚，则齐其庶几乎（大概差不多了吧）。（《孟子二章·庄暴见孟子》）

③览物之情，得无异乎（能够没有什么不同吗）？（《岳阳楼记》）

④成反复自念，得无教我猎虫所耶（莫不是指教我捉蟋蟀的处所吧）？（《促织》）

（7）不亦······乎：相当于“不是······吗”、“不也······吗”。例如：

①求剑若此，不亦惑乎（不是糊涂吗）？（《刻舟求剑》）

②学而时习之，不亦说乎（不也是一种快乐吗？）（《论语》六则）

（8）何以······为：相当于“用（要）······做什么”。例如：

①匈奴未灭，何以家为？（匈奴没有消灭掉，要家做什么呢？）（《汉书》）

②然则又何以兵为？（既然这样，那么用兵做什么？）（《荀子·议兵》）

（9）何······为：相当于“为什么······呢”。例如：

①如今人方为刀俎，我为鱼肉，何辞为（为什么要告辞呢）？（《鸿门宴》）

②天之亡我，我何渡为（我为什么要渡过江呢）？（《史记·项羽本纪》）

③吾英王，奚跪汝为（我是英王，为什么要跪着你呢）？（《陈玉成》）

（10）有以、无以：相当于“有可以拿来······的”、“没有可以拿来······的”。例如：

①项王未有以应。（项王没有可以拿来回答的话。）（《鸿门宴》）

②吾终当有以活汝。（我总会有可以拿来使你活下去的办法。）（《中山狼传》）

③不积小流无以成江海。（不汇集细小的水流，就没有可以拿来汇成江海的东西。）（《劝学》）

④河曲智叟亡（通“无”）以应。（河曲智叟没有可以拿来回答的话。）（《愚公移山》）

（11）有所、无所：相当于“有什么”，“没有什么”。例如：

①恐太后玉体之有所郄也。（担心太后贵体有什么不舒适。）（《触龙说赵太后》）

②顾念蓄劣物终无所用，不如拼搏一笑。（可是心想养着这低劣的东西终究没有什么用，还不如让它拼一拼，换得一笑。）（《促织》）

（12）何所······、谁······者：相当于“所······（者）何”、“······者谁”。例如：

①问女何所思，问女何所忆。（问女思想的是什么，问女怀念的是什么。）（《木兰诗》）

②谁可使者？（可以出使秦国的是谁？）（《廉颇蔺相如列传》）
\newpage